\documentclass[11pt]{article}
\usepackage[T1]{fontenc}
\usepackage[paperwidth=175mm,paperheight=210mm,margin=14mm]{geometry}
\usepackage{lmodern,amsmath,amssymb,microtype}
\pagestyle{empty}
\begin{document}
\begin{center}
{\Large\bfseries Split divisor boxes\par}
\vspace{2mm}
{\large An exact ES forcing criterion}\par
\vspace{2mm}
The Clankers\quad---\quad13 September 2026
\end{center}
\hrule\vspace{4mm}
Retain the entire prime-exponent box $B_e=\prod_i[-e_i,e_i]_{\mathbb Z}$
modulo $R=4a-p$. Split its budget $e=f+b$, and pull the original targets
through the removed part:
\[
\mathcal T=\{-1,-p,-p^{-1}\},\qquad
\mathcal F_b=\mathcal T\,v(B_b)^{-1}.
\]
\[
\boxed{v(B_e)\cap\mathcal T\ne\varnothing
\iff v(B_f)\cap\mathcal F_b\ne\varnothing.}
\]
In each quadratic-character cell, exact integer/parity minimization gives
\[
\boxed{\text{energy}<\text{target-avoidance minimum}
\ \Longrightarrow\ \text{an original ES witness}.}
\]
\textbf{Strict example.} At $p=1201$, $a=312$, $R=47$, remove one
power of $2$ from the exponent budget. The unsplit test does not certify.
The partial negative-character cell has
\[
\text{mass}=30,\qquad\text{energy}=46<50=\text{minimum}.
\]
Its marked inverse returns
\[
\boxed{\frac4{1201}=\frac1{312}+\frac1{9608}+\frac1{46839}.}
\]
\textbf{Finite comparison.} Among 9531 specified shells, the earlier test
certifies 240; adding the split tests certifies 438. Fourteen additions use
strict finite minima. These are shell counts, not new prime coverage.

\vspace{3mm}\hrule\vspace{3mm}
\small
The original exponent multiplicities, labels, signs and inverse fibres remain
explicit. Failed states retain the exact denominator defect and affine inverse.
The three-page note supplies proofs and reproducible standard-library checks.
No universal ES proof, priority claim or independent review is asserted.
This card summarizes the proof, rather than replacing it.
\end{document}
